\documentclass[11pt,a4paper]{article}

\usepackage[T1]{fontenc}
\usepackage[utf8]{inputenc}
\usepackage[english]{babel}
\usepackage[margin=1in]{geometry}
\usepackage{microtype}
\usepackage{hyperref}
\usepackage[backend=biber,style=numeric,sorting=none]{biblatex}
\addbibresource{references.bib}

\begin{document}

\begin{titlepage}
\centering
\vspace*{0.18\textheight}
{\LARGE\bfseries An Ablation Study of Auxiliary Inputs for Lightweight Neural Upscaling in Real-Time Rendering\par}
\vspace{2.5cm}
{\large Zhanbolot Kydyraliev \quad and \quad Amin Soltobekov\par}
\vspace{0.5cm}
{\large Affiliation to be confirmed\par}
\vfill
{\large Working manuscript --- \today\par}
\end{titlepage}

\begin{abstract}
This manuscript is in preparation. The abstract will be written after the research question, experimental protocol, and results are finalized.
\end{abstract}

\section{Introduction}
Neural upscaling reconstructs a higher-resolution image from lower-resolution rendering data. A real-time renderer can also provide auxiliary signals, such as depth and motion vectors, that may help reconstruction. Each additional input has a potential quality benefit and a cost in data preparation, memory, and runtime. This study is planned as a controlled ablation of selected auxiliary inputs for a lightweight upscaler. The exact model, inputs, hardware, and evaluation conditions will be fixed in the experiment protocol before results are collected.

We address five research questions about the individual and combined
contributions of depth and surface normals, the effect of motion-aligned
history, the measured costs of these variants, their performance on unseen
scenes, and where in an image the inputs help or hurt. The questions and
study design are given in Section~\ref{sec:rq}.

% TODO: Add a source-backed motivation and define lightweight/real-time.
% TODO: State the final research question and contributions after scope agreement.

\section{Background and Related Work}
% TODO: Synthesize verified sources by method, auxiliary input, and evaluation setting.
% Cite a verified entry from references.bib with \cite{CitationKey}.
The following work motivates the use of temporal history, rendering buffers,
and sample reconstruction in the proposed study.

\subsection{Temporal Anti-Aliasing and Supersampling}
Low-resolution rendering can produce visible aliasing because rendered pixels
sample the scene at discrete positions~\cite{xiao2020neural}. Karis's
temporal anti-aliasing method for Unreal Engine 4 combines samples from
successive frames after shifting the sampling positions between frames
\cite{karis2014high}. Motion vectors reproject the previous result into the
current frame. Inaccurate motion vectors can smear the result, while stale
history can cause ghosting; Karis uses neighborhood clipping to limit the
influence of incompatible history. These observations motivate our comparison
between unaligned and motion-aligned history. That comparison tests the
effect of alignment within our proposed upscaler, not the full temporal
anti-aliasing method.

\subsection{Neural Supersampling for Real-Time Rendering}
Xiao et al.~\cite{xiao2020neural} present neural supersampling for real-time
rendering. Their method reconstructs a high-resolution frame from
low-resolution rendered color and uses depth, motion vectors, and information
from previous frames. Motion vectors align earlier frames with the current
frame, and the method learns to use the aligned history where it is useful.
Their work demonstrates a system that combines current-frame rendering buffers
with temporal information. Our proposed study asks a narrower question: how
quality and measured cost change across selected input variants under a
shared experimental protocol. In particular, V4 and V5 compare history
without and with motion alignment; they do not measure the effect of motion
vectors independently of temporal history.

\subsection{Pixels as Samples}
Smith~\cite{smith1995pixel} explains that a pixel represents a sample of an
underlying image rather than a small square of uniform color. Upscaling can
therefore be viewed as reconstruction from samples. This perspective makes
classical interpolation, such as bicubic interpolation, a useful candidate
baseline for a learned method. Separately, the proposed study will render
low-resolution inputs natively at their intended input resolution, rather
than create them by
filtered downsampling of high-resolution frames. This choice is intended to
retain the aliasing produced by the selected rendering process; it is a
decision about our experimental data, not a claim made by Smith.

\section{Method}
\subsection{Research Questions and Study Design}
\label{sec:rq}
We study how auxiliary rendering inputs affect a lightweight neural
upscaler. We ask five research questions.
\begin{description}
\item[RQ1 (individual contribution)] Relative to color-only input, how do
depth and surface normals each affect reconstruction quality, and how does
motion-aligned history compare with unaligned history?
\item[RQ2 (combination)] Do the quality gains from depth and normals overlap
in the single-frame setting, and what is the incremental effect of adding
them together to motion-aligned history?
\item[RQ3 (cost)] How do the selected runtime and memory measurements vary
across the variants on the provisional target GPU,
relative to their reconstruction quality?
\item[RQ4 (generalization)] Do quality differences between variants persist
on held-out scenes, and how do they vary between those scenes?
\item[RQ5 (where the inputs help)] Where in the image do auxiliary inputs
help or hurt, for example at geometric edges, on thin structures, and in
regions newly revealed by motion?
\end{description}

\paragraph{Hypotheses.}
Before running experiments, we state expectations that the study can confirm
or refute.
\begin{description}
\item[H1] Adding depth improves quality over color-only input.
\item[H2] Adding surface normals improves quality over color-only input.
\item[H3] Depth and normals are partly redundant, so using both improves
quality by less than the sum of their individual gains.
\item[H4] History warped with motion vectors improves quality over history
used without alignment. Incorrect motion can smear temporal results
\cite{karis2014high}, while motion-based alignment is used in neural
supersampling~\cite{xiao2020neural}.
\item[H5] Improvements of selected auxiliary-input variants over their
matched baselines persist on scene-disjoint test scenes.
\end{description}

\paragraph{Study design.}
We plan to use a shared base architecture while varying the inputs. A shared
design does not imply identical parameter counts: changing the number of
input channels changes the first layer, and temporal variants may require
additional processing. The final architecture and treatment of temporal
history will be fixed before experimentation. Training data, loss,
optimizer, and schedule should be held constant across variants where
applicable. We plan to use multiple random seeds; their number remains to be
set. Parameter counts and computational complexity will be reported so that
quality differences are not attributed solely to the inputs if model
capacity also differs. Possible controls include matching parameter budgets
or adding a capacity-matched comparison, subject to protocol finalization.

Table~\ref{tab:variants} lists the core variants. V1 versus V0 isolates
depth, and V2 versus V0 isolates normals. V3 versus V1 estimates the
incremental effect of normals with depth, while V3 versus V2 estimates the
incremental effect of depth with normals; together with V0, these
comparisons test overlap in the single-frame setting. V4 versus V0
estimates the effect of adding unaligned history, while V5 versus V4 tests
motion alignment of that history, provided the temporal model and history
inputs are otherwise held constant. V6 versus V5 estimates the combined
addition of depth and normals
to aligned history; V6 versus V3 compares the combined temporal variant with
its single-frame counterpart. These comparisons do not independently
identify every depth, normal, and temporal interaction. In particular,
the matrix does not isolate motion vectors without history or the individual
effects of depth and normals when aligned history is present. Variants that
add only depth or only normals to aligned history could be considered as
optional extensions if the study scope permits.

Training, validation, and test scenes will be disjoint at the scene level;
frames from a test scene will not be used for model fitting or protocol
selection. Validation scenes will support model and protocol choices, while
held-out test scenes will support the reported comparisons. Scenes with a
defined distribution shift would require a separate out-of-distribution
test; being unseen during training alone does not establish such a shift.
RQ1 and RQ2 use the variant comparisons above; RQ3 compares quality with
the selected cost measurements; RQ4 compares variants across held-out
scenes; and RQ5 examines spatial error patterns, including error maps.
Candidate metrics and cost dimensions are discussed in
Section~\ref{sec:metrics}. The temporal variants are a planned extension,
and their implementation and the final protocol remain to be fixed.

\begin{table}[h]
\centering
\caption{Planned core input variants (V0 is the color-only baseline).}
\label{tab:variants}
\begin{tabular}{llcccc}
\hline
ID & Group & Color & Depth & Normals & History \\
\hline
V0 & Single frame & yes & -- & -- & -- \\
V1 & Single frame & yes & yes & -- & -- \\
V2 & Single frame & yes & -- & yes & -- \\
V3 & Single frame & yes & yes & yes & -- \\
V4 & Temporal & yes & -- & -- & unaligned \\
V5 & Temporal & yes & -- & -- & motion aligned \\
V6 & Temporal & yes & yes & yes & motion aligned \\
\hline
\end{tabular}
\end{table}
% TODO: fix the number of seeds, the scenes, and the train/test split in
% Section 3.2 and 3.3 before collecting results.

\subsection{Model, Inputs, and Data}
% TODO: Specify architecture, input preprocessing, scenes, splits, resolution, and scale.
The architecture, input preprocessing, scenes, resolution, and upscaling
factor are to be selected. The training, validation, and test scene lists
and the method for producing aligned or unaligned history will be fixed
before any experiments are run.

\subsection{Evaluation Protocol}
\label{sec:metrics}
% TODO: Specify baselines, quality metrics, latency/memory measurement, and hardware.
Candidate image-quality measures include PSNR and SSIM; a perceptual metric
may be added if it suits the final reference images. These measures would
support RQ1, RQ2, and RQ4. A temporal consistency measure may be useful for
RQ5 and the temporal comparisons, provided it can distinguish instability
from genuine scene motion. Error maps and selected image regions can also
support RQ5. The final measures, aggregation across frames and scenes, and
statistical treatment across seeds will be specified before testing.

For RQ3, possible cost measures include neural inference latency and GPU
memory consumption. Auxiliary-buffer preparation, history warping and
storage, data transfers, and total upscaling-pipeline time may also matter
for deployment. The protocol must specify which of these are measured, how
their boundaries are defined, and the target hardware and timing procedure.
Any costs outside those boundaries will be identified as unmeasured, so a
quality--cost comparison is not mistaken for a complete pipeline benchmark.
The measurement protocol and any interpolation baseline remain to be fixed
before the experiments.

\section{Results}
% TODO: Add tables, figures, run IDs, and uncertainty only after experiments.
Results have not yet been collected.

\subsection{Contribution of Individual and Combined Inputs (RQ1, RQ2)}
% TODO: results not yet collected.

\subsection{Runtime and Memory Cost (RQ3)}
% TODO: results not yet collected.

\subsection{Generalization to Unseen Scenes (RQ4)}
% TODO: results not yet collected.

\subsection{Where the Inputs Help or Hurt (RQ5)}
% TODO: results not yet collected.

\section{Discussion}
% TODO: Interpret effects and interactions, compare with cited work, and discuss limits.
This section will answer RQ1 to RQ5 in turn, state whether hypotheses H1 to
H5 were supported by the results, and discuss the limitations of the study.

\section{Conclusion}
% TODO: Answer the final research question using measured evidence.
Conclusions will be written after the analysis.

\clearpage
\section*{References}
\addcontentsline{toc}{section}{References}
\printbibliography[heading=none]

\end{document}
